This study presented an efficient and robust pipeline for the automated classification of ISUP grades in prostate histopathology images, combining ordinal regression, Focal Loss, and entropy-based filtering with variants of the EfficientNet architecture. The results demonstrated that the proposed approach substantially improves the clinical coherence of predictions, reduces severe ordinal errors, and outperforms individual models by leveraging ensembles based on probabilistic averaging.

The simple averaging ensemble achieved the best overall performance, with QWK = 0.879 and an accuracy of 69.8\%, highlighting that the combination of complementary architectures provides greater stability and generalization. Beyond quantitative performance, the model exhibited more consistent behavior from a clinical perspective, with errors predominantly concentrated between adjacent classes.

The main limitation is the use of a single dataset, albeit multicentric and heterogeneous. Future evaluations on independent external datasets are necessary to confirm the method's robustness and generalization. Furthermore, the greater difficulty observed in the intermediate grades suggests that specific strategies for these classes may represent a promising path for further improvements.
Future work includes investigating Transformer-based architectures, especially Vision Transformers and hybrid CNN-Transformer models, as well as incorporating active learning methods and automatic uncertainty-based curation strategies, aiming to make the system even more robust and clinically applicable.